\chapter{UN Sustainable Development Goals}
\section{Mapping with UN Sustainable Development Goals}
\paragraph{}The 2030 Agenda for Sustainable Development, adopted by 193 Member States at the UN General Assembly Summit in September 2015, defines 17 Sustainable Development Goals (SDGs) covering economic, social, and environmental dimensions of global development. Every project undertaken in academic and professional contexts should be evaluated for its contribution toward one or more of these goals.

\paragraph{}SMPTMS --- Smart Property Tax Management System --- directly aligns with \textbf{SDG 11: Sustainable Cities and Communities}, which aims to make cities inclusive, safe, resilient, and sustainable. A key enabler of sustainable urban governance is the financial independence of municipalities. Property tax is the most significant own-source revenue stream for Urban Local Bodies in India. By digitizing and streamlining property tax collection for Jaipur Nagar Nigam, SMPTMS contributes to stronger municipal finances, more transparent governance, and improved delivery of urban services for all citizens.

\begin{table}[h]
    \begin{tabular}{|l|p{6cm}|p{8cm}|}
        \hline
        \multicolumn{3}{|l|}{\textbf{SDG Title: SDG 11 -- Sustainable Cities and Communities}} \\
        \hline
        \textbf{S.No} & \textbf{Objective} & \textbf{Outcome} \\
        \hline
        1 & Improve municipal revenue collection through digitization of property tax administration & Increased tax compliance leading to better funding for civic infrastructure including roads, sanitation, and water supply in Jaipur \\
        \hline
        2 & Make property tax assessment transparent and accessible to all citizens regardless of digital literacy level & Reduced scope for corruption and manual manipulation in tax calculations through a publicly visible, formula-based system \\
        \hline
        3 & Enable online payment to remove physical barriers for citizens in remote wards who cannot easily visit municipal offices & Higher payment rates and reduced administrative burden on ward-level staff, enabling faster revenue realization \\
        \hline
        4 & Provide GIS-based property mapping to support evidence-based urban planning and property record maintenance & Accurate, spatially-referenced property data enabling smarter infrastructure investment decisions by Jaipur Nagar Nigam \\
        \hline
    \end{tabular}
    \caption{SDG 11 Mapping for SMPTMS}
    \label{tab:sdg}
\end{table}
